\documentclass[11pt]{article}
\usepackage[a4paper,margin=27mm]{geometry}
\usepackage[T1]{fontenc}
\usepackage{lmodern,amsmath,amssymb,amsthm,microtype}
\usepackage[colorlinks=true,linkcolor=blue!50!black,citecolor=blue!50!black,urlcolor=blue!50!black]{hyperref}
\usepackage{xcolor}
\hypersetup{pdftitle={Every word, unequal frequencies: digits of power Lambert series},pdfauthor={},pdfsubject={Disjunctivity, positive word frequencies, and nonnormality of power Lambert series}}
\newcommand{\T}{\mathbb T}
\newcommand{\N}{\mathbb N}
\newcommand{\Z}{\mathbb Z}
\newcommand{\E}{\mathbb E}
\newcommand{\ind}{\mathbf 1}
\newcommand{\dd}{\,d}
\newtheorem{theorem}{Theorem}
\newtheorem{lemma}{Lemma}
\setlength{\parindent}{0pt}
\setlength{\parskip}{5pt plus 1pt minus 1pt}
\setlength{\emergencystretch}{1.5em}
\raggedbottom

\begin{document}
\begin{center}
{\Large\bfseries Every word, unequal frequencies:\\[3pt]
digits of power Lambert series}
\end{center}
\vspace{3pt}

For integers $b,r\ge2$, consider
\[
L_{b,r}=\sum_{m=1}^{\infty}\frac{1}{b^{m^r}-1}.
\]
A real number is \emph{disjunctive in base $b$} if every finite word over
$\{0,\ldots,b-1\}$ occurs in its fractional base-$b$ expansion. It is
\emph{normal in base $b$} if every word of length $\ell$ has limiting
frequency $b^{-\ell}$. Occurrences here may overlap, and frequencies are
taken along all initial segments.

\begin{theorem}\label{thm:main}
For every $b,r\ge2$, every finite base-$b$ word has a strictly positive
limiting frequency in $L_{b,r}$, but $L_{b,r}$ is not normal in base $b$.
In particular, it is disjunctive and irrational.

More precisely, the empirical measures of $\{b^nL_{b,r}\}$ converge to an
explicit probability measure $\nu_{b,r}$ on $\T=\mathbb R/\Z$.
This measure is invariant and ergodic under $T_b(x)=bx\pmod1$, has no
atoms, gives positive mass to every nonempty open set, and has
Kolmogorov--Sinai entropy zero.
\end{theorem}

The proof has two parts. Congruences force a prescribed word while an
average estimate controls the carry from later terms. Separately, a
periodic approximation identifies the limiting distribution and shows
that its entropy is zero. The arithmetic construction develops the
Chowla--Erd\H{o}s method as used by Duverney, Duverney--Tachiya, and
Vandehey~\cite{Duverney,DT,Vandehey}. Irrationality of these constants
already follows from \cite[Corollary 1.2]{DT}; the issue here is their
digit distribution.

\section{Forcing a word}

Fix $b,r\ge2$ and write $L=L_{b,r}$. For $v\ge1$, let
\[
c_r(v)=\#\{d\ge1:d^r\mid v\}
      =\prod_{p^a\parallel v}\left(\left\lfloor\frac a r\right\rfloor+1\right).
\]
Expanding each denominator as a geometric series gives
\begin{equation}\label{eq:carry}
L=\sum_{v\ge1}c_r(v)b^{-v},\qquad
\{b^nL\}=\left\{\sum_{j\ge1}c_r(n+j)b^{-j}\right\}.
\end{equation}
All rearrangements are justified by nonnegativity and $c_r(v)\le v$.
The coefficients need not be digits: this is why the carry matters.

\subsection*{Two elementary estimates}

\begin{lemma}[Means on a progression]\label{lem:mean}
For positive integers $u,Q$,
\[
\lim_{M\to\infty}\frac1M\sum_{m=0}^{M-1}c_r(u+Qm)
=\sum_{d\ge1}\frac{\gcd(d^r,Q)}{d^r}
  \ind_{\gcd(d^r,Q)\mid u}.
\]
In the Euler product on the right, primes not dividing $Q$ contribute
factors $(1-p^{-r})^{-1}$. If $p^s\parallel Q$ with
$s=r(B-1)+1$ for an integer $B\ge1$, its factor is at most $B+1$;
it is $1$ if $p\nmid u$.
\end{lemma}

\begin{proof}
The congruence $d^r\mid u+Qm$ is either insoluble or one residue class
modulo $d^r/\gcd(d^r,Q)$. This gives the stated limit for each $d$.
The average of its indicator is at most $(u+Q)/d^r$: the progression
values are distinct positive integers bounded by $u+QM$.
Since $\sum d^{-r}<\infty$, dominated convergence permits summation.

Multiplicativity gives the Euler product. If $p^s\nmid u$, the valuation
at $p$ is fixed below $s$, and its factor is at most $B$.
If $p^s\mid u$, the factor is
\[
B+\sum_{k\ge1}p^{s-r(B-1+k)}
 =B+\frac{p}{p^r-1}\le B+1.
\]
If $p\nmid u$, the valuation is zero and the factor is $1$.
\end{proof}

The same domination allows passage through a weighted tail: for fixed
$u,Q,H$, the mean of $\sum_{j>H}c_r(u+Qm+j)b^{-j}$ is the sum of its
termwise means. Indeed, the summable majorant is
$(u+Q+j)\zeta(r)b^{-j}$.

\begin{lemma}[Power-free values]\label{lem:sieve}
If $u,v\ge1$ and $\gcd(u,v)=1$, then the indices $m\ge0$ for which
$u+vm$ is $r$-free have lower density at least $1/4$.
Here $r$-free means divisible by no $p^r$ with $p$ prime.
\end{lemma}

\begin{proof}
A prime dividing $v$ divides no progression value. Every other prime
excludes at most $M/p^r+1$ of the first $M$ indices. Only primes at most
$\sqrt{u+vM}$ can exclude a value, since $r\ge2$. Thus the excluded
proportion is at most
\[
\sum_p p^{-r}+O_{u,v}(M^{-1/2})
\le\sum_{n=2}^{\infty}n^{-2}+O_{u,v}(M^{-1/2})
\le\frac34+O_{u,v}(M^{-1/2}).
\]
The last inequality follows from $1/4+\int_2^\infty x^{-2}\dd x=3/4$.
\end{proof}

\subsection*{The congruences and the carry}

Let a prescribed word have length $\ell\ge1$ and integer value
$0\le w<b^\ell$. Its digit interval is
$I_w=[w/b^\ell,(w+1)/b^\ell)$. Put
\[
t=\ell+1,\qquad A=bw+1,\qquad e=r(A-1),\qquad
\eta=\frac{1}{2b^t}.
\]
The closed interval
\begin{equation}\label{eq:guard}
K_w=\left[\frac A{b^t},\frac A{b^t}+\eta\right]
\quad\hbox{lies strictly inside }I_w.
\end{equation}
The extra digit $1$ leaves room for a small positive carry.

Choose a large integer $H\ge\max(t,3)$; its size will be fixed below.
For each $j\in\{0,\ldots,H\}\setminus\{t\}$ choose $H$ primes, all
distinct and greater than $H^3+H+2$. Call these the primes assigned to
$j$, let $D$ be their total number, and put
\[
s=r(b-1)+1,\qquad Q=2^{e+1}\prod_{p\text{ assigned}}p^s.
\]
We will use only $D\le H(H+1)$. The Chinese remainder theorem gives
$0\le a<Q$ satisfying
\[
a+t\equiv2^e\pmod{2^{e+1}},\qquad
a+j\equiv p^{s-1}\pmod{p^s}
\quad\text{for each prime $p$ assigned to $j$}.
\]
Take $n=a+Qm$ with $m\ge1$. Each prime assigned to $j$ has exact
valuation $r(b-1)$ in $n+j$, so contributes a factor $b$ to $c_r(n+j)$.
Consequently,
\begin{equation}\label{eq:integral}
b^j\mid c_r(n+j)\qquad(1\le j\le H,\ j\ne t).
\end{equation}

Write $n+t=2^e u_m$. The $u_m$ form an odd positive arithmetic
progression of step $Q/2^e$. No assigned prime divides $n+t$: a prime
assigned to $j$ would then divide $t-j$, although $0<|t-j|\le H<p$.
Thus the cofactor progression is coprime. By Lemma~\ref{lem:sieve},
$u_m$ is $r$-free on a set of lower density at least $1/4$.
At those indices,
\begin{equation}\label{eq:target}
c_r(n+t)=\left(\frac e r+1\right)c_r(u_m)=A.
\end{equation}

It remains to control $R_H(n)=\sum_{j>H}c_r(n+j)b^{-j}$.
Let $M_j$ be the mean of $c_r(a+Qm+j)$ on this progression, and put
$C=2(A+1)$. If $H<j\le H^3$, no assigned prime divides $n+j$:
subtracting its assigned shift gives a nonzero difference smaller than
the prime. Lemma~\ref{lem:mean}, together with $\zeta(r)\le2$, gives
\[
M_j\le C\quad(H<j\le H^3),\qquad
M_j\le C(b+1)^{H(H+1)}\quad(j>H^3).
\]
The factor at $2$ is at most $A+1$; each assigned prime costs at most
$b+1$ when it is not excluded. Summing the two geometric tails yields
\begin{equation}\label{eq:tailmean}
\lim_{M\to\infty}\frac1M\sum_{m=1}^{M}R_H(a+Qm)
\le \frac C{b-1}\left(b^{-H}+(b+1)^{H(H+1)}b^{-H^3}\right).
\end{equation}
For $H\ge3$, the second term in parentheses is at most $b^{-H}$:
use $b+1\le b^2$ and $2H(H+1)+H\le H^3$.
Choose $H$ so that the right side of~\eqref{eq:tailmean} is less than
$\eta/16$, and only then choose the primes. The estimate is independent
of their sizes.

Markov's inequality shows that $R_H(n)\ge\eta$ on a set of upper
density less than $1/16$ within the progression. Removing these indices
from the power-free set leaves positive lower density. On the remaining
indices, \eqref{eq:carry}, \eqref{eq:integral}, and \eqref{eq:target} give
\[
\{b^nL\}=\frac A{b^t}+R_H(n)\in K_w.
\]
A set of positive lower density in one fixed progression has positive
lower density among all integers. We have therefore proved:
\begin{equation}\label{eq:visits}
\liminf_{N\to\infty}\frac1N
 \#\{0\le n<N:\{b^nL\}\in K_w\}>0
\quad\text{for every word }w.
\end{equation}
In particular $L$ is disjunctive. It is irrational, since the radix orbit
of a rational number is finite, whereas \eqref{eq:visits} makes this
orbit dense in the circle.

\section{The limiting distribution}

The second part of the proof starts from a different expression for the
same carry. For $d\ge1$, put
\[
g_d(n)=\frac{b^{n\bmod d}}{b^d-1},\qquad
W(n)=\sum_{m\ge1}g_{m^r}(n),\qquad
W_U(n)=\sum_{m=1}^{U}g_{m^r}(n).
\]
Summing the geometric series over $j$ with $d\mid n+j$ shows that
$W(n)=\sum_{j\ge1}c_r(n+j)b^{-j}$. Thus $W(n)\pmod1=\{b^nL\}$.

\subsection*{An approximation valid for every prefix}

The $d$ values of $g_d$ in a period increase from its minimum to its
maximum and have total $1/(b-1)$. The average over an initial part of
that period is therefore at most the complete-period average.
Splitting any prefix into complete periods and one initial part gives
\[
\frac1N\sum_{n=0}^{N-1}g_d(n)\le\frac{1}{d(b-1)}
\qquad(N\ge1).
\]
Summing over $m>U$ proves the key estimate
\begin{equation}\label{eq:uniform}
\sup_{N\ge1}\frac1N\sum_{n=0}^{N-1}|W(n)-W_U(n)|
\le\varepsilon_U:=\frac{1}{b-1}\sum_{m>U}m^{-r}
\longrightarrow0.
\end{equation}
Each $W_U$ is periodic, with period dividing $(U!)^r$.
The uniformity in $N$ lets us pass from these periodic sequences to the
actual digit sequence.

\subsection*{A compact group model}

Record all power residues at once:
\[
G=\overline{\{(n\bmod m^r)_{m\ge1}:n\in\Z\}}
\ \subseteq\ \prod_{m\ge1}\Z/m^r\Z.
\]
This is a compact group. Let $\mu$ be its Haar probability measure and
$R(x)=x+\mathbf1$, where $\mathbf1=(1\bmod m^r)_{m\ge1}$.
Any finite collection of coordinates ranges over a single finite cyclic
orbit. Haar measure projects to the uniform measure on that orbit.
The rotation $R$ is ergodic: an invariant $L^2$ function is invariant
under the dense subgroup generated by $\mathbf1$, hence, by continuity
of translations in $L^2$, under every translation, and is constant
almost everywhere. This is the usual compact-group rotation argument
\cite{MF}. Related residue models in arithmetic dynamics appear in
\cite{CS,CV}.

For $x\in G$, let $\rho_m(x)\in\{0,\ldots,m^r-1\}$ represent its
$m$th coordinate. Define
\[
W(x)=\sum_{m\ge1}\frac{b^{\rho_m(x)}}{b^{m^r}-1},\qquad
\Phi(x)=W(x)\pmod1,\qquad \nu=\Phi_*\mu.
\]
Every coordinate is uniform, so
\[
\int_G W\dd\mu=\frac{\zeta(r)}{b-1}<\infty.
\]
The series is therefore finite almost everywhere; define $\Phi=0$ on
its null exceptional set. Each summand satisfies
\[
b\frac{b^{\rho_m(x)}}{b^{m^r}-1}
-\frac{b^{\rho_m(Rx)}}{b^{m^r}-1}\in\Z.
\]
Reduce finite partial sums modulo $1$ and pass to the limit where both
series converge. This proves
\begin{equation}\label{eq:factor}
\Phi\circ R=T_b\circ\Phi\quad\mu\text{-almost everywhere}.
\end{equation}
Thus $\nu$ is invariant and ergodic under $T_b$.

Let $W_U(x)$ be the corresponding finite sum on $G$. Its Haar tail
satisfies $\int|W-W_U|\dd\mu=\varepsilon_U$.
If $f$ is a Lipschitz function on the circle, \eqref{eq:uniform} bounds
the difference between the integer averages of $f(W)$ and $f(W_U)$ by
$\operatorname{Lip}(f)\varepsilon_U$. The same bound holds for their
Haar integrals. For fixed $U$, the integer average of $f(W_U)$ converges
to its Haar integral, because it runs periodically through the finite
cyclic quotient. Hence
\[
\limsup_{N\to\infty}
\left|\frac1N\sum_{n<N}f(\{b^nL\})-\int f\dd\nu\right|
\le2\operatorname{Lip}(f)\varepsilon_U.
\]
Letting $U\to\infty$, and using uniform approximation of continuous
circle functions by Lipschitz functions, proves
\begin{equation}\label{eq:weak}
\frac1N\sum_{n<N}\delta_{\{b^nL\}}\ \Longrightarrow\ \nu.
\end{equation}
This identifies the distribution of the specified orbit, not just of
almost every point of the model.

\subsection*{From the distribution to word frequencies}

For the closed interval $K_w$, weak convergence gives
\[
\nu(K_w)\ge\limsup_{N\to\infty}\frac1N
\#\{n<N:\{b^nL\}\in K_w\}>0
\]
by \eqref{eq:visits}. Every nonempty open arc contains a digit interval,
so $\nu$ has full support.

An ergodic invariant probability measure with an atom is supported on a
finite periodic orbit. Indeed, atom masses cannot decrease under forward
iteration, so the orbit of an atom must eventually repeat. The resulting
cycle has positive measure and is invariant modulo a null set;
ergodicity gives it full measure. Full support on the circle excludes
this possibility. Thus $\nu$ has no atoms.

The boundary of $I_w$ is finite and therefore has measure zero.
Equation~\eqref{eq:weak} now gives the limiting word frequency
\[
\lim_{N\to\infty}\frac1N
\#\{n<N:\{b^nL\}\in I_w\}=\nu(I_w)>0.
\]

\section{Zero entropy and failure of normality}

Here entropy measures the average information per symbol in a stationary
sequence of finite labels. For a finite random variable $X$ with
probabilities $p_i$, its Shannon entropy is
$H(X)=-\sum_i p_i\log p_i$, with $0\log0=0$.
We prove directly that every finite measurable labeling of the group
rotation has entropy rate zero.

Let $X:G\to\{1,\ldots,k\}$ be such a labeling. Sets depending on
finitely many coordinates generate the Borel sigma algebra of $G$.
Consequently $X$ can be approximated by a finite-coordinate labeling
$Y$ with mismatch probability $p=\mu(X\ne Y)$ as small as desired.
For completeness, measurable sets approximable in measure by the
finite-coordinate algebra form a sigma algebra; approximate the finitely
many fibers of $X$ and assign any overlaps in a fixed order.
There is an integer $P\ge1$ such that $Y\circ R^P=Y$.

Define the error label $E$ to be $*$ when $X=Y$, and $X$ otherwise.
The first $n$ labels of $X$ along an orbit are determined by the first
$P$ labels of $Y$ and the first $n$ error labels. By the elementary
reconstruction and subadditivity inequalities for Shannon entropy,
and invariance of $\mu$,
\[
H(X,X\circ R,\ldots,X\circ R^{n-1})
\le P\log k+nH(E).
\]
Moreover,
\[
H(E)\le -p\log p-(1-p)\log(1-p)+p\log k\longrightarrow0
\quad(p\to0).
\]
Divide by $n$, let $n\to\infty$, and then improve the approximation.
Every finite labeling has entropy rate zero, so their supremum, the
Kolmogorov--Sinai entropy of $R$, is zero. Iterating
\eqref{eq:factor} outside a null set shows that pulling a circle
labeling back through $\Phi$ preserves all its block distributions.
Therefore $T_b$ also has entropy zero with respect to $\nu$.

If $L$ were normal, its limiting word frequencies would give each of
the $b^n$ length-$n$ digit intervals measure $b^{-n}$ under $\nu$.
The digit labeling would then have block entropy $n\log b$, and hence
positive entropy rate $\log b$. This contradicts the preceding result
and completes the proof of Theorem~\ref{thm:main}.

\paragraph{Lean verification.}
The accompanying source package proves all conclusions of
Theorem~\ref{thm:main} in
\texttt{PowerLambert.powerLambert\_full\_theorem}, with only the
hypotheses $b,r\ge2$. It uses Lean 4.34.0-rc2, Mathlib~\cite{Mathlib},
and the Shannon entropy library from the Polynomial Freiman--Ruzsa
formalization~\cite{PFR}; exact revisions and build instructions are
recorded in the package. The axiom audit contains only Lean's standard
\texttt{propext}, \texttt{Classical.choice}, and \texttt{Quot.sound}.
The formalization also proves nonnormality directly from periodic
approximation; the paper uses the shorter entropy consequence above.

\begin{thebibliography}{9}
\small\raggedright
\setlength{\itemsep}{3pt}
\bibitem{Duverney}
D. Duverney, \emph{Arithmetical functions and irrationality of Lambert
series}, AIP Conference Proceedings \textbf{1385} (2011), 5--16.
\href{https://danielduverney.fr/documents/theorie-des-nombres/TokyoNT.pdf}{Author manuscript}.

\bibitem{DT}
D. Duverney and Y. Tachiya, \emph{Refinement of the Chowla--Erd\H{o}s
method and linear independence of certain Lambert series},
Forum Mathematicum \textbf{31} (2019), 1557--1566.
\href{https://danielduverney.fr/documents/theorie-des-nombres/DuverneyTachiya190522.pdf}{Author manuscript}.

\bibitem{Vandehey}
J. Vandehey, \emph{On an incomplete argument of Erd\H{o}s on the
irrationality of Lambert series}, Integers \textbf{13} (2013), A58.
\href{https://arxiv.org/abs/1206.0340}{arXiv:1206.0340}.

\bibitem{MF}
M. Mirzakhani, lectures; T. Feng, notes, \emph{Introduction to Ergodic
Theory} (2014).
\href{https://math.berkeley.edu/~fengt/ergodic_theory.pdf}{Lecture notes}.

\bibitem{CS}
F. Cellarosi and Ya. G. Sinai, \emph{Ergodic properties of square-free
numbers}, Journal of the European Mathematical Society \textbf{15}
(2013), 1343--1374.
\href{https://ems.press/journals/jems/articles/11437}{Journal article}.

\bibitem{CV}
F. Cellarosi and I. Vinogradov, \emph{Ergodic properties of $k$-free
integers in number fields}, Journal of Modern Dynamics \textbf{7}
(2013), 461--488.
\href{https://arxiv.org/abs/1304.0214}{arXiv:1304.0214}.

\bibitem{Mathlib}
The Mathlib contributors, \emph{The Lean mathematical library},
revision \texttt{de2ef68216c6}.
\href{https://github.com/leanprover-community/mathlib4/tree/de2ef68216c6074f338c8e61890ee0a379ddfb9b}{Source repository}.

\bibitem{PFR}
The Polynomial Freiman--Ruzsa formalization contributors,
\emph{Polynomial Freiman--Ruzsa conjecture formalization},
Shannon entropy library, revision \texttt{3d7898164ebf}.
\href{https://github.com/teorth/pfr/tree/3d7898164ebff70a809dce618f9082a7b39e7850}{Source repository}.
\end{thebibliography}

\medskip
{\footnotesize Text licensed under
\href{https://creativecommons.org/licenses/by/4.0/}{CC BY 4.0}.
The accompanying Lean code is licensed under Apache 2.0.}
\end{document}
